\documentclass[uplatex,nomag*,a4paper,11pt]{jsarticle}
\usepackage[dvipdfmx,margin=25mm]{geometry}
\usepackage{tabularx}
\begin{document}
\section{風見輪の刻み}
薄い紙の円盤に同じ間隔で印を付けた。\\
最初の印を上に向ける。
\subsection{二枚の比較}
\noindent\begin{tabularx}{\linewidth}{p{4zw}XX}
刻み & 一枚目の形 & 二枚目の形 \\
外周 & 丸い島 & 細い岬 \\
内周 & 短い橋 & 深い入り江
\end{tabularx}
\paragraph{記録}白い印の数を書き留める。
\vspace*{164mm}
\section{裏面の印}
\begin{figure}[t]
\centering\rule{28mm}{4mm}
\caption{印を写すための帯}
\end{figure}
\noindent\begin{minipage}[t]{.46\linewidth}
\subsection*{外輪}円盤を返して外周をなぞる。\\
一番深い刻みに白い点を打つ。
\end{minipage}\hfill
\begin{minipage}[t]{.46\linewidth}
\subsection*{内輪}内周の印を紙へ移す。\\
最後の印を青で囲む。
\end{minipage}
\end{document}
